% --- Archon physics formalization source begin ---
% archon:physics
% archon:covers PhyXMiniProblems/problem_phyx_mini_0567.lean
% archon:source-report reports/phyx_mini/problem_phyx_mini_0567.source.json

\chapter{Physics problem phyx\_mini\_0567}
\label{ch:PhyXMiniProblems_problem_phyx_mini_0567}

\paragraph{Problem source.}
\#\# Physical scenario

The light from a nearby star is observed to be shifted toward blue by 2\%, i.e. \$\textbackslash{}frac\{\textbackslash{}Delta \textbackslash{}lambda\}\{\textbackslash{}lambda\} = -0.02\$.

\#\# Question

How fast is it moving?

\#\# Answer choices

- A: A: 0.55c

- B: B: 0.0198c

- C: C: 0.95c

- D: D: 10.2c

\#\# Figure caption (auxiliary; use the image as primary evidence)

The image shows two diagrams with coordinate systems, representing different perspectives involving a muon and Earth.

**Left Diagram:**
- A coordinate system labeled \textbackslash{}(O\_\textbackslash{}mu\textbackslash{}) with axes \textbackslash{}(x\_\textbackslash{}mu\textbackslash{}) (horizontal) and \textbackslash{}(y\_\textbackslash{}mu\textbackslash{}) (vertical). 
- A red circle labeled "μ" is located on the negative \textbackslash{}(y\_\textbackslash{}mu\textbackslash{}) axis.
- An arrow labeled "V" points downwards from the circle.
- A red stick figure stands at the origin \textbackslash{}(O\_e\textbackslash{}) on the positive \textbackslash{}(x\_e\textbackslash{}) axis, with Earth represented as a horizontal black line.
- Text: \textbackslash{}(y\_e \textbackslash{}equiv -y\_\textbackslash{}mu\textbackslash{})

**Right Diagram:**
- A coordinate system labeled \textbackslash{}(O\_e\textbackslash{}) with axes \textbackslash{}(x\_e\textbackslash{}) (horizontal) and \textbackslash{}(y\_e\textbackslash{}) (vertical).
- Earth is depicted above the \textbackslash{}(x\_e\textbackslash{}) axis as a horizontal black line with squiggly edges below, labeled "earth."
- An arrow labeled "V" points downwards from the Earth line.
- A red stick figure is situated at \textbackslash{}(O\_\textbackslash{}mu\textbackslash{}) on the positive \textbackslash{}(x\_\textbackslash{}mu\textbackslash{}) axis, with a red circle labeled "μ" on the \textbackslash{}(x\_\textbackslash{}mu\textbackslash{}) axis near the origin.
- Text: \textbackslash{}(y\_\textbackslash{}mu \textbackslash{}equiv -y\_e\textbackslash{})

These diagrams illustrate reference frames involving motion and spatial perception differences between two observers.

\#\# Dataset metadata

Relativity; Physical Model Grounding Reasoning, Numerical Reasoning

\paragraph{Recorded answer/context.}
B: B: 0.0198c

\paragraph{Figure/image path.}
/root/proposal\_for\_physic/science-mango/phyx\_mini\_run/phyx\_data/test\_image/567.png

\paragraph{Formalization target.}
create a compiling Lean file with sorry bodies at `PhyXMiniProblems/problem\_phyx\_mini\_0567.lean`. The Lean declarations must preserve the physical quantities, dimensions or dimensional roles, figure labels, governing-law hypotheses, and final relation expressed by this problem.
Use Mathlib/Physlib names found through LeanExplore where available. If a physics API is missing, introduce faithful local abstractions rather than scalar placeholder aliases.

\begin{theorem}[Physics formalization target]
\label{thm:physics:phyx_mini_0567:target}
\lean{PhyXMiniProblems.ProblemPhyXMini0567.problem_phyx_mini_0567}
\uses{def:physics:phyx-mini-0567:phyxminiproblems-problemphyxmini0567-speedreadout, def:physics:phyx-mini-0567:phyxminiproblems-problemphyxmini0567-vacuumlightspeedreadout, def:physics:phyx-mini-0567:phyxminiproblems-problemphyxmini0567-matchessuppliedauxiliaryfigure, def:physics:phyx-mini-0567:phyxminiproblems-problemphyxmini0567-stellarblueshiftsetup, def:physics:phyx-mini-0567:phyxminiproblems-problemphyxmini0567-speedfractionoflight, def:physics:phyx-mini-0567:phyxminiproblems-problemphyxmini0567-matchesstellarblueshiftscenario, def:physics:phyx-mini-0567:phyxminiproblems-problemphyxmini0567-matchestwopercentwavelengthblueshift, def:physics:phyx-mini-0567:phyxminiproblems-problemphyxmini0567-hasphysicalstellardopplerparameters, def:physics:phyx-mini-0567:phyxminiproblems-problemphyxmini0567-satisfiesapproachingsourcerelativisticdopplerlaw, def:physics:phyx-mini-0567:phyxminiproblems-problemphyxmini0567-answerchoice, def:physics:phyx-mini-0567:phyxminiproblems-problemphyxmini0567-isclosestdisplayedspeedchoice}
The assigned autoformalize agent should translate this physics problem into Lean declarations in the covered file, with theorem and lemma proof bodies written as `by sorry`.
\end{theorem}
\begin{proof}
This is an autoformalization task, not a proof task. Produce statements that can later be proved by the physics prover without weakening the physical model.
\end{proof}

% --- Archon declaration topology begin ---
\section{Declaration topology}
\begin{definition}[Wavelength Quantity]
  \label{def:physics:phyx-mini-0567:phyxminiproblems-problemphyxmini0567-wavelengthquantity}
  \lean{PhyXMiniProblems.ProblemPhyXMini0567.WavelengthQuantity}
  A nonnegative physical wavelength, carrying length dimension
\end{definition}
\begin{proof}
  This is the declared model definition of Wavelength Quantity.
\end{proof}
\begin{definition}[Signed Wavelength Change Quantity]
  \label{def:physics:phyx-mini-0567:phyxminiproblems-problemphyxmini0567-signedwavelengthchangequantity}
  \lean{PhyXMiniProblems.ProblemPhyXMini0567.SignedWavelengthChangeQuantity}
  A signed physical wavelength change Δλ
\end{definition}
\begin{proof}
  This is the declared model definition of Signed Wavelength Change Quantity.
\end{proof}
\begin{definition}[Speed Quantity]
  \label{def:physics:phyx-mini-0567:phyxminiproblems-problemphyxmini0567-speedquantity}
  \lean{PhyXMiniProblems.ProblemPhyXMini0567.SpeedQuantity}
  A nonnegative physical radial speed
\end{definition}
\begin{proof}
  This is the declared model definition of Speed Quantity.
\end{proof}
\begin{definition}[wavelength Readout]
  \label{def:physics:phyx-mini-0567:phyxminiproblems-problemphyxmini0567-wavelengthreadout}
  \lean{PhyXMiniProblems.ProblemPhyXMini0567.wavelengthReadout}
  \uses{def:physics:phyx-mini-0567:phyxminiproblems-problemphyxmini0567-wavelengthquantity}
  Read a wavelength in a selected unit of length
\end{definition}
\begin{proof}
  This is the declared model definition of wavelength Readout.
\end{proof}
\begin{definition}[signed Wavelength Change Readout]
  \label{def:physics:phyx-mini-0567:phyxminiproblems-problemphyxmini0567-signedwavelengthchangereadout}
  \lean{PhyXMiniProblems.ProblemPhyXMini0567.signedWavelengthChangeReadout}
  \uses{def:physics:phyx-mini-0567:phyxminiproblems-problemphyxmini0567-signedwavelengthchangequantity}
  Read the signed wavelength change in a selected unit of length
\end{definition}
\begin{proof}
  This is the declared model definition of signed Wavelength Change Readout.
\end{proof}
\begin{definition}[speed Readout]
  \label{def:physics:phyx-mini-0567:phyxminiproblems-problemphyxmini0567-speedreadout}
  \lean{PhyXMiniProblems.ProblemPhyXMini0567.speedReadout}
  \uses{def:physics:phyx-mini-0567:phyxminiproblems-problemphyxmini0567-speedquantity}
  Read a speed in a selected length unit per selected time unit
\end{definition}
\begin{proof}
  This is the declared model definition of speed Readout.
\end{proof}
\begin{definition}[vacuum Light Speed Readout]
  \label{def:physics:phyx-mini-0567:phyxminiproblems-problemphyxmini0567-vacuumlightspeedreadout}
  \lean{PhyXMiniProblems.ProblemPhyXMini0567.vacuumLightSpeedReadout}
  Read Physlib's dimensionful vacuum speed of light in compatible units
\end{definition}
\begin{proof}
  This is the declared model definition of vacuum Light Speed Readout.
\end{proof}
\begin{definition}[Astronomical Source Kind]
  \label{def:physics:phyx-mini-0567:phyxminiproblems-problemphyxmini0567-astronomicalsourcekind}
  \lean{PhyXMiniProblems.ProblemPhyXMini0567.AstronomicalSourceKind}
  The astronomical source classification stated in the problem
\end{definition}
\begin{proof}
  This is the declared model definition of Astronomical Source Kind.
\end{proof}
\begin{definition}[Observer Location]
  \label{def:physics:phyx-mini-0567:phyxminiproblems-problemphyxmini0567-observerlocation}
  \lean{PhyXMiniProblems.ProblemPhyXMini0567.ObserverLocation}
  The observing location relative to which the radial motion is measured
\end{definition}
\begin{proof}
  This is the declared model definition of Observer Location.
\end{proof}
\begin{definition}[Light Signal Kind]
  \label{def:physics:phyx-mini-0567:phyxminiproblems-problemphyxmini0567-lightsignalkind}
  \lean{PhyXMiniProblems.ProblemPhyXMini0567.LightSignalKind}
  The light observable whose displacement is measured spectroscopically
\end{definition}
\begin{proof}
  This is the declared model definition of Light Signal Kind.
\end{proof}
\begin{definition}[Doppler Frame]
  \label{def:physics:phyx-mini-0567:phyxminiproblems-problemphyxmini0567-dopplerframe}
  \lean{PhyXMiniProblems.ProblemPhyXMini0567.DopplerFrame}
  The two inertial frames used for emission and observation
\end{definition}
\begin{proof}
  This is the declared model definition of Doppler Frame.
\end{proof}
\begin{definition}[Radial Motion]
  \label{def:physics:phyx-mini-0567:phyxminiproblems-problemphyxmini0567-radialmotion}
  \lean{PhyXMiniProblems.ProblemPhyXMini0567.RadialMotion}
  Direction of the star's radial motion in the observer's convention
\end{definition}
\begin{proof}
  This is the declared model definition of Radial Motion.
\end{proof}
\begin{definition}[Auxiliary Figure Panel]
  \label{def:physics:phyx-mini-0567:phyxminiproblems-problemphyxmini0567-auxiliaryfigurepanel}
  \lean{PhyXMiniProblems.ProblemPhyXMini0567.AuxiliaryFigurePanel}
  The left and right panels in the supplied auxiliary image
\end{definition}
\begin{proof}
  This is the declared model definition of Auxiliary Figure Panel.
\end{proof}
\begin{definition}[Auxiliary Reference Frame]
  \label{def:physics:phyx-mini-0567:phyxminiproblems-problemphyxmini0567-auxiliaryreferenceframe}
  \lean{PhyXMiniProblems.ProblemPhyXMini0567.AuxiliaryReferenceFrame}
  The two reference frames labelled by the image origins O\_mu and O\_e
\end{definition}
\begin{proof}
  This is the declared model definition of Auxiliary Reference Frame.
\end{proof}
\begin{definition}[Auxiliary Axis Label]
  \label{def:physics:phyx-mini-0567:phyxminiproblems-problemphyxmini0567-auxiliaryaxislabel}
  \lean{PhyXMiniProblems.ProblemPhyXMini0567.AuxiliaryAxisLabel}
  Literal coordinate-axis labels visible in the image
\end{definition}
\begin{proof}
  This is the declared model definition of Auxiliary Axis Label.
\end{proof}
\begin{definition}[Auxiliary Origin Label]
  \label{def:physics:phyx-mini-0567:phyxminiproblems-problemphyxmini0567-auxiliaryoriginlabel}
  \lean{PhyXMiniProblems.ProblemPhyXMini0567.AuxiliaryOriginLabel}
  Origin labels visible in the image
\end{definition}
\begin{proof}
  This is the declared model definition of Auxiliary Origin Label.
\end{proof}
\begin{definition}[Auxiliary Figure Object]
  \label{def:physics:phyx-mini-0567:phyxminiproblems-problemphyxmini0567-auxiliaryfigureobject}
  \lean{PhyXMiniProblems.ProblemPhyXMini0567.AuxiliaryFigureObject}
  Physical objects drawn in both panels
\end{definition}
\begin{proof}
  This is the declared model definition of Auxiliary Figure Object.
\end{proof}
\begin{definition}[Auxiliary Object Glyph]
  \label{def:physics:phyx-mini-0567:phyxminiproblems-problemphyxmini0567-auxiliaryobjectglyph}
  \lean{PhyXMiniProblems.ProblemPhyXMini0567.AuxiliaryObjectGlyph}
  Glyph styles used to distinguish the three objects in the raster
\end{definition}
\begin{proof}
  This is the declared model definition of Auxiliary Object Glyph.
\end{proof}
\begin{definition}[Vertical Direction]
  \label{def:physics:phyx-mini-0567:phyxminiproblems-problemphyxmini0567-verticaldirection}
  \lean{PhyXMiniProblems.ProblemPhyXMini0567.VerticalDirection}
  Vertical direction of each arrow labelled V
\end{definition}
\begin{proof}
  This is the declared model definition of Vertical Direction.
\end{proof}
\begin{definition}[Vertical Coordinate Relation]
  \label{def:physics:phyx-mini-0567:phyxminiproblems-problemphyxmini0567-verticalcoordinaterelation}
  \lean{PhyXMiniProblems.ProblemPhyXMini0567.VerticalCoordinateRelation}
  The opposite-axis convention printed at the top of either panel
\end{definition}
\begin{proof}
  This is the declared model definition of Vertical Coordinate Relation.
\end{proof}
\begin{definition}[Muon Earth Reference Frame Figure]
  \label{def:physics:phyx-mini-0567:phyxminiproblems-problemphyxmini0567-muonearthreferenceframefigure}
  \lean{PhyXMiniProblems.ProblemPhyXMini0567.MuonEarthReferenceFrameFigure}
  \uses{def:physics:phyx-mini-0567:phyxminiproblems-problemphyxmini0567-auxiliaryfigurepanel, def:physics:phyx-mini-0567:phyxminiproblems-problemphyxmini0567-auxiliaryreferenceframe, def:physics:phyx-mini-0567:phyxminiproblems-problemphyxmini0567-auxiliaryaxislabel, def:physics:phyx-mini-0567:phyxminiproblems-problemphyxmini0567-auxiliaryoriginlabel, def:physics:phyx-mini-0567:phyxminiproblems-problemphyxmini0567-auxiliaryfigureobject, def:physics:phyx-mini-0567:phyxminiproblems-problemphyxmini0567-auxiliaryobjectglyph, def:physics:phyx-mini-0567:phyxminiproblems-problemphyxmini0567-verticaldirection, def:physics:phyx-mini-0567:phyxminiproblems-problemphyxmini0567-verticalcoordinaterelation}
  Qualitative and literal information transcribed from image 567. The image has no quantitative spatial scale and no stellar wavelength or speed label
\end{definition}
\begin{proof}
  This is the declared model definition of Muon Earth Reference Frame Figure.
\end{proof}
\begin{definition}[Matches Supplied Auxiliary Figure]
  \label{def:physics:phyx-mini-0567:phyxminiproblems-problemphyxmini0567-matchessuppliedauxiliaryfigure}
  \lean{PhyXMiniProblems.ProblemPhyXMini0567.MatchesSuppliedAuxiliaryFigure}
  \uses{def:physics:phyx-mini-0567:phyxminiproblems-problemphyxmini0567-muonearthreferenceframefigure}
  Primary-image evidence. The left panel is organized around O\_mu, with x\_mu, y\_mu, and x\_e shown and the muon's V arrow downward. The right panel is organized around O\_e, with x\_e, y\_e, and x\_mu shown and the Earth's V arrow downward. The displayed coordinate conventions are y\_e ≡ -y\_mu and y\_mu ≡ -y\_e respectively
\end{definition}
\begin{proof}
  This is the declared model definition of Matches Supplied Auxiliary Figure.
\end{proof}
\begin{definition}[Stellar Blueshift Setup]
  \label{def:physics:phyx-mini-0567:phyxminiproblems-problemphyxmini0567-stellarblueshiftsetup}
  \lean{PhyXMiniProblems.ProblemPhyXMini0567.StellarBlueshiftSetup}
  \uses{def:physics:phyx-mini-0567:phyxminiproblems-problemphyxmini0567-wavelengthquantity, def:physics:phyx-mini-0567:phyxminiproblems-problemphyxmini0567-signedwavelengthchangequantity, def:physics:phyx-mini-0567:phyxminiproblems-problemphyxmini0567-speedquantity, def:physics:phyx-mini-0567:phyxminiproblems-problemphyxmini0567-astronomicalsourcekind, def:physics:phyx-mini-0567:phyxminiproblems-problemphyxmini0567-observerlocation, def:physics:phyx-mini-0567:phyxminiproblems-problemphyxmini0567-lightsignalkind, def:physics:phyx-mini-0567:phyxminiproblems-problemphyxmini0567-dopplerframe, def:physics:phyx-mini-0567:phyxminiproblems-problemphyxmini0567-radialmotion, def:physics:phyx-mini-0567:phyxminiproblems-problemphyxmini0567-muonearthreferenceframefigure}
  The physical quantities and role assignments of the observation. starRadialSpeedRelativeToObserver is independent data constrained only by the Doppler law below; it is not defined from the recorded answer choice
\end{definition}
\begin{proof}
  This is the declared model definition of Stellar Blueshift Setup.
\end{proof}
\begin{definition}[speed Fraction Of Light]
  \label{def:physics:phyx-mini-0567:phyxminiproblems-problemphyxmini0567-speedfractionoflight}
  \lean{PhyXMiniProblems.ProblemPhyXMini0567.speedFractionOfLight}
  \uses{def:physics:phyx-mini-0567:phyxminiproblems-problemphyxmini0567-speedreadout, def:physics:phyx-mini-0567:phyxminiproblems-problemphyxmini0567-vacuumlightspeedreadout, def:physics:phyx-mini-0567:phyxminiproblems-problemphyxmini0567-stellarblueshiftsetup}
  The dimensionless speed ratio β = v/c, derived from physical speeds
\end{definition}
\begin{proof}
  This is the declared model definition of speed Fraction Of Light.
\end{proof}
\begin{definition}[Matches Stellar Blueshift Scenario]
  \label{def:physics:phyx-mini-0567:phyxminiproblems-problemphyxmini0567-matchesstellarblueshiftscenario}
  \lean{PhyXMiniProblems.ProblemPhyXMini0567.MatchesStellarBlueshiftScenario}
  \uses{def:physics:phyx-mini-0567:phyxminiproblems-problemphyxmini0567-stellarblueshiftsetup}
  The nearby stellar source, Earth observer, frame assignments, and approaching radial direction stated or implied by a blueshift
\end{definition}
\begin{proof}
  This is the declared model definition of Matches Stellar Blueshift Scenario.
\end{proof}
\begin{definition}[Matches Two Percent Wavelength Blueshift]
  \label{def:physics:phyx-mini-0567:phyxminiproblems-problemphyxmini0567-matchestwopercentwavelengthblueshift}
  \lean{PhyXMiniProblems.ProblemPhyXMini0567.MatchesTwoPercentWavelengthBlueshift}
  \uses{def:physics:phyx-mini-0567:phyxminiproblems-problemphyxmini0567-wavelengthreadout, def:physics:phyx-mini-0567:phyxminiproblems-problemphyxmini0567-signedwavelengthchangereadout, def:physics:phyx-mini-0567:phyxminiproblems-problemphyxmini0567-stellarblueshiftsetup}
  The literal problem readout Δλ / λ = -0.02. Here Δλ = λ\_observed - λ\_emitted, and the denominator is the emitted rest wavelength. The equations are imposed in every length unit to make their dimensionless meaning explicit. No speed value occurs in these premises
\end{definition}
\begin{proof}
  This is the declared model definition of Matches Two Percent Wavelength Blueshift.
\end{proof}
\begin{definition}[Has Physical Stellar Doppler Parameters]
  \label{def:physics:phyx-mini-0567:phyxminiproblems-problemphyxmini0567-hasphysicalstellardopplerparameters}
  \lean{PhyXMiniProblems.ProblemPhyXMini0567.HasPhysicalStellarDopplerParameters}
  \uses{def:physics:phyx-mini-0567:phyxminiproblems-problemphyxmini0567-wavelengthreadout, def:physics:phyx-mini-0567:phyxminiproblems-problemphyxmini0567-vacuumlightspeedreadout, def:physics:phyx-mini-0567:phyxminiproblems-problemphyxmini0567-stellarblueshiftsetup, def:physics:phyx-mini-0567:phyxminiproblems-problemphyxmini0567-speedfractionoflight}
  Positivity and the ordinary subluminal branch of the stellar observation. These conditions select the physical branch of the square root but do not assign a numerical value to the star's speed
\end{definition}
\begin{proof}
  This is the declared model definition of Has Physical Stellar Doppler Parameters.
\end{proof}
\begin{definition}[Satisfies Approaching Source Relativistic Doppler Law]
  \label{def:physics:phyx-mini-0567:phyxminiproblems-problemphyxmini0567-satisfiesapproachingsourcerelativisticdopplerlaw}
  \lean{PhyXMiniProblems.ProblemPhyXMini0567.SatisfiesApproachingSourceRelativisticDopplerLaw}
  \uses{def:physics:phyx-mini-0567:phyxminiproblems-problemphyxmini0567-wavelengthreadout, def:physics:phyx-mini-0567:phyxminiproblems-problemphyxmini0567-stellarblueshiftsetup, def:physics:phyx-mini-0567:phyxminiproblems-problemphyxmini0567-speedfractionoflight}
  The longitudinal relativistic Doppler law for an approaching source, λ\_observed / λ\_emitted = sqrt ((1 - β) / (1 + β)). This is a general governing relation between the independent wavelengths and speed. It contains neither the problem-specific exact speed nor an answer choice
\end{definition}
\begin{proof}
  This is the declared model definition of Satisfies Approaching Source Relativistic Doppler Law.
\end{proof}
\begin{definition}[Answer Choice]
  \label{def:physics:phyx-mini-0567:phyxminiproblems-problemphyxmini0567-answerchoice}
  \lean{PhyXMiniProblems.ProblemPhyXMini0567.AnswerChoice}
  Labels of the four speed choices printed with the problem
\end{definition}
\begin{proof}
  This is the declared model definition of Answer Choice.
\end{proof}
\begin{definition}[displayed Speed Fraction Of Light]
  \label{def:physics:phyx-mini-0567:phyxminiproblems-problemphyxmini0567-displayedspeedfractionoflight}
  \lean{PhyXMiniProblems.ProblemPhyXMini0567.displayedSpeedFractionOfLight}
  \uses{def:physics:phyx-mini-0567:phyxminiproblems-problemphyxmini0567-answerchoice}
  The nonnegative coefficient of c printed beside each answer label
\end{definition}
\begin{proof}
  This is the declared model definition of displayed Speed Fraction Of Light.
\end{proof}
\begin{definition}[recorded Dataset Answer]
  \label{def:physics:phyx-mini-0567:phyxminiproblems-problemphyxmini0567-recordeddatasetanswer}
  \lean{PhyXMiniProblems.ProblemPhyXMini0567.recordedDatasetAnswer}
  \uses{def:physics:phyx-mini-0567:phyxminiproblems-problemphyxmini0567-answerchoice}
  The answer label recorded by the source dataset; it is not a premise
\end{definition}
\begin{proof}
  This is the declared model definition of recorded Dataset Answer.
\end{proof}
\begin{definition}[Is Closest Displayed Speed Choice]
  \label{def:physics:phyx-mini-0567:phyxminiproblems-problemphyxmini0567-isclosestdisplayedspeedchoice}
  \lean{PhyXMiniProblems.ProblemPhyXMini0567.IsClosestDisplayedSpeedChoice}
  \uses{def:physics:phyx-mini-0567:phyxminiproblems-problemphyxmini0567-stellarblueshiftsetup, def:physics:phyx-mini-0567:phyxminiproblems-problemphyxmini0567-speedfractionoflight, def:physics:phyx-mini-0567:phyxminiproblems-problemphyxmini0567-answerchoice, def:physics:phyx-mini-0567:phyxminiproblems-problemphyxmini0567-displayedspeedfractionoflight}
  A displayed choice is selected by ordinary multiple-choice proximity to the exact modeled speed fraction. This matters because the printed 0.0198 c is close to, but not equal to, the exact consequence of the written wavelength shift
\end{definition}
\begin{proof}
  This is the declared model definition of Is Closest Displayed Speed Choice.
\end{proof}
% --- Archon declaration topology end ---
% --- Archon physics formalization source end ---
